%! Convolutional Neural Nets — Unit 8, Lesson 8.2 — Guided Notes (Answer Key)
\documentclass[10pt]{article}
\usepackage{linalg-article}
\usepackage{linalg-key}   % pulls in -boxes; do NOT also load -boxes
\newcommand{\TallMath}[1]{$\displaystyle #1\rule[-1.4em]{0pt}{3.2em}$}
\newcommand{\vv}{\mathbf{v}}
\newcommand{\bb}{\mathbf{b}}
\newcommand{\ww}{\mathbf{w}}
\newcommand{\zero}{\mathbf{0}}
\newcommand{\relu}{\mathrm{ReLU}}

\begin{document}

\pageheader{Unit 8, Lesson 8.2}{Guided Notes --- Answer Key}
\namedateperiod

\vspace{0.12in}

\begin{objectivebox}
\textbf{By the end of this lesson, I will be able to\dots}
\begin{itemize}\setlength{\itemsep}{2pt}
  \item slide a small \emph{filter} across a signal, taking a weighted sum (dot product) at each spot, and read where the filter \emph{fires};
  \item write a convolution as a matrix times a vector whose weights \emph{repeat down the diagonals} (a shared-weight matrix), and count how many weights it saves;
  \item explain how a convolutional layer $\relu(A\vv+\bb)$ reuses one small filter everywhere --- so it needs far fewer weights and detects the same pattern \emph{wherever} it appears.
\end{itemize}
\end{objectivebox}

\vspace{0.06in}

\begin{vocabbox}
\small
Fill in each term as we build it together.
\vspace{2pt}
\vocabans{Filter (kernel) --- the small set of weights we slide}{a short list of weights (e.g.\ $[-1,1]$) that acts as a pattern-detector; the same filter is used at every position}
\vocabans{Convolution --- slide the filter, dot product at each spot}{line the filter up with each window of the signal and take the weighted sum; the results form the output}
\vocabans{Weight sharing --- the same filter weights reused everywhere}{one small filter is applied at every location, so a few weights do the whole job instead of a separate weight per connection}
\vocabans{Convolution matrix --- weights repeat down the diagonals}{the matrix form of a convolution: the filter's numbers appear in every row, shifted over by one (a constant-diagonal matrix)}
\vocabans{Feature map --- the output signal, high where the pattern is}{the result of the convolution; a large value marks a spot where the filter's pattern is present}
\end{vocabbox}

\begin{hookbox}
Last lesson a \textbf{layer} $\relu(A\vv+\bb)$ used a full weight matrix $A$ --- a separate weight for every
connection. That is fine for a few numbers, but a photo has \emph{millions} of pixels.
\begin{itemize}\setlength{\itemsep}{1pt}
  \item A ``bright edge'' looks the same whether it sits in the top-left corner or the middle of a picture. Should a network learn a \emph{separate} weight to spot that edge at every location? \ansline{No --- that would waste enormous numbers of weights relearning the same pattern. Better to learn \emph{one} edge-detector and reuse it everywhere.}
  \item If one small pattern-detector could be \emph{reused} at every spot, what would that do to the number of weights we have to learn? \ansline{It shrinks it dramatically --- a handful of shared weights instead of one weight per pixel-pair, no matter how big the image is.}
\end{itemize}
\end{hookbox}

\begin{notesbox}{1. Too Many Weights --- and the Pattern Repeats}
In a full layer, the matrix $A$ has one weight for \emph{every} input-output connection. For an image with a
million pixels that is around a \emph{trillion} weights --- far too many to learn, and wasteful: a useful
pattern (an edge, a corner) is the \textbf{same shape} wherever it appears. The fix is one idea:
\begin{center}
\emph{use a \textbf{small} set of weights, and \textbf{reuse} the same weights at every position.}
\end{center}
That small set of weights is called a \textbf{filter} (or \textbf{kernel}). Sliding it across the input to
look for its pattern everywhere is called a \textbf{convolution}.
\end{notesbox}

\begin{notesbox}{2. Convolution: Slide a Small Filter}
Take the two-weight filter $\ww=[-1,\,1]$ and a signal $\vv=[\,2,\,2,\,2,\,6,\,6\,]$ (flat, then a jump up).
\emph{Slide} the filter along the signal; at each spot line it up with a window and take the dot product
(here $\ww\cdot[v_i,v_{i+1}]=-v_i+v_{i+1}$, the ``jump'' from one value to the next):
\[
  y_1={\color{keyred}\mathbf{0}},\qquad y_2={\color{keyred}\mathbf{0}},\qquad y_3={\color{keyred}\mathbf{4}},\qquad y_4={\color{keyred}\mathbf{0}}.
\]
(At the window $[2,2]$: $-2+2=0$. At $[2,6]$: $-2+6=4$.) The output is $[\,0,\,0,\,4,\,0\,]$: it is $0$ on the
flat parts and \emph{fires} (nonzero) exactly at the \textbf{edge} where $2$ jumps to $6$. One filter, slid
everywhere, is an \textbf{edge detector}. The output signal is called a \textbf{feature map} --- big where
the pattern is present.

\begin{center}
\begin{tikzpicture}[scale=0.62]
  % LEFT: the input signal with the filter window sliding over the edge
  \begin{scope}[shift={(-3.4,0)}]
    \draw[->, charcoal, thin] (-0.3,0)--(5.6,0) node[right]{\scriptsize position};
    \draw[->, charcoal, thin] (0,-0.4)--(0,3.4) node[above]{\scriptsize value};
    % step signal 2,2,2,6,6 at x=1..5, scale y by 0.42
    \draw[ultra thick, royalblue]
      (1,0.84)--(2,0.84)--(3,0.84)--(4,2.52)--(5,2.52);
    \foreach \x/\y in {1/0.84,2/0.84,3/0.84,4/2.52,5/2.52}
      \fill[royalblue] (\x,\y) circle (1.7pt);
    % sliding window over the edge (positions 3-4)
    \draw[goldacc, very thick, dashed] (2.75,0.5)--(4.25,0.5)--(4.25,2.9)--(2.75,2.9)--cycle;
    \node[goldacc] at (3.5,3.3) {\scriptsize filter $[-1,1]$ slides};
    \node[charcoal] at (2.7,-0.75) {\scriptsize signal $[2,2,2,6,6]$};
  \end{scope}
  % ARROW
  \draw[->, very thick, goldacc] (2.55,1.3)--(4.35,1.3) node[midway, above]{\scriptsize convolve};
  % RIGHT: the feature map, spike at the edge
  \begin{scope}[shift={(5.7,0)}]
    \draw[->, charcoal, thin] (-0.3,0)--(5.0,0) node[right]{\scriptsize position};
    \draw[->, charcoal, thin] (0,-0.4)--(0,3.4) node[above]{\scriptsize output};
    % bars 0,0,4,0 at x=1..4, scale 0.42
    \foreach \x/\y in {1/0,2/0,3/1.68,4/0}{
      \draw[ultra thick, burgundy] (\x,0)--(\x,\y);
      \fill[burgundy] (\x,\y) circle (1.7pt);
    }
    \node[burgundy] at (3.0,2.15) {\scriptsize $4$};
    \node[charcoal] at (3.0,-0.75) {\scriptsize fires at the edge};
  \end{scope}
\end{tikzpicture}
\end{center}
\end{notesbox}

\begin{notesbox}{3. It Is a Matrix --- with Shared Weights}
Sliding a filter is just a matrix times a vector (Unit~1). Put the filter weights $-1,1$ into each row,
shifted one step to the right each time:
\[
  A=\begin{bmatrix} -1 & 1 & 0 & 0 & 0\\ 0 & -1 & 1 & 0 & 0\\ 0 & 0 & -1 & 1 & 0\\ 0 & 0 & 0 & -1 & 1 \end{bmatrix},
  \qquad A\vv = A\begin{bmatrix}2\\2\\2\\6\\6\end{bmatrix}=\begin{bmatrix}{\color{keyred}\mathbf{0}}\\[-1pt] {\color{keyred}\mathbf{0}}\\[-1pt] {\color{keyred}\mathbf{4}}\\[-1pt] {\color{keyred}\mathbf{0}}\end{bmatrix}.
\]
The \emph{same} two numbers repeat down the diagonals: this is a \textbf{convolution matrix} (a
constant-diagonal matrix). \textbf{Weight sharing} is what makes it special.

\vspace{2pt}
\textbf{Counting the savings.} A \emph{full} $4\times5$ matrix would have $4\cdot 5=20$ separate weights. This
convolution has only \ans{$2$}. For a real $1000\times1000$ image a full layer needs about $10^{12}$
weights; a small $3\times3$ filter needs just $9$ --- reused everywhere.
\end{notesbox}

\begin{notesbox}{4. A Convolutional Layer --- and Why It Matters}
A convolutional layer is the \emph{same} layer as before, $\relu(A\vv+\bb)$ --- only now $A$ is a
convolution matrix (shared, sliding weights) and $\bb$ is a bias. On our feature map $[0,0,4,0]$ with a bias
$-1$:
\[
  \relu\!\big([0,0,4,0]+(-1)\big)=\relu([-1,-1,3,-1])=\big[\ {\color{keyred}\mathbf{0}},\ {\color{keyred}\mathbf{0}},\ {\color{keyred}\mathbf{3}},\ {\color{keyred}\mathbf{0}}\ \big].
\]
The ReLU (with a bias \emph{threshold}) keeps only the strong response --- ``the edge is really here'' --- and
zeroes the rest.

\vspace{2pt}
\textbf{Why it matters.} Because the one filter is reused at \emph{every} position, it detects its pattern
\emph{wherever} that pattern sits --- move the edge and the same filter still fires (just at the new spot).
That is why these networks recognize a cat in the corner or the center of a photo. And weight sharing cuts a
trillion weights down to a handful, so the network is small enough to actually \emph{learn} (Lesson~8.3).

\vspace{2pt}
\textbf{The road ahead.} \textbf{8.3} minimizing loss by gradient descent\; $\bullet$\; \textbf{8.4} mean,
variance, and covariance.
\end{notesbox}

\begin{practicebox}
Show your work.
\begin{enumerate}[leftmargin=1.6em, itemsep=4pt]
  \item Slide the filter $\ww=[-1,1]$ (output $=-v_i+v_{i+1}$) across $\vv=[4,4,4,9]$. What is the output, and at which spot does it fire? \ansline{$y_1=4-4=0$, $y_2=4-4=0$, $y_3=9-4=5$. Output $[0,0,5]$; it fires at the third spot --- the edge where $4$ jumps to $9$.}
  \item A full layer for a signal of length $6$ (six outputs) would need $36$ weights. How many weights does a length-$3$ filter, reused everywhere, need instead? \ansline{Just $3$ --- the three filter weights are shared across all positions, no matter how long the signal is.}
  \item Take the feature map from problem~1 and apply $\relu(\,\cdot\,+\bb)$ with bias $\bb=-2$. What comes out, and what does the leftover value mean? \ansline{$\relu([0,0,5]+(-2))=\relu([-2,-2,3])=[0,0,3]$. Only the strong edge survives; the leftover $3$ says the pattern is genuinely present there, while the $-2$ threshold clears away the flat spots.}
\end{enumerate}
\end{practicebox}

\begin{teachernote}
The spine is Section~2 (slide $\ww=[-1,1]$ across $[2,2,2,6,6]$ to get the feature map $[0,0,4,0]$ --- a
\emph{spike at the edge}) and Section~3 (the same slide written as a constant-diagonal matrix $A\vv$, with
only $2$ shared weights instead of $20$). Section~1 sets up the motivation --- full layers have too many
weights and relearn the same pattern everywhere --- and Section~4 lands the payoff: a conv layer is still
$\relu(A\vv+\bb)$, but weight sharing makes it small \emph{and} translation-invariant (it detects a pattern
wherever it sits). Keep numbers tiny and lean on the picture: flat signal $\to$ zeros, a jump $\to$ a spike.
Most common slips: sign/order errors in $-v_i+v_{i+1}$ (emphasize ``next minus current''); expecting a full
matrix instead of noticing the repeated diagonals; and thinking the filter must change from spot to spot (the
whole point is that it does \emph{not}). Sentence to land: one small filter, reused everywhere, finds its
pattern anywhere --- with almost no weights.
\end{teachernote}

\end{document}
